\begin{theorem}\label{night:P1.1}
\StatementRecord{night:P1.1}{proved}{}{Definitions and the displayed constructions}
The finite homology monodromies have orders $5,4$, respectively, and
characteristic polynomials
$(T-1)^2\Phi _5(T)$ and $(T-1)^2(T+1)^2(T^2+1)$.
For $N=M_0-I$ we have $N^2=0$, $\operatorname{rank}N=3$, and
nonzero Smith factors $1,1,1$.  Moreover
\[
 \Lambda^{\langle M_1,M_2\rangle}=0,\qquad
 \Lambda_{\langle M_1,M_2\rangle}=\mathbb Z,
\]
where the quotient homomorphism is $\ell(x)=x_0$.
On the cohomology lattice the corresponding groups are $\mathbb Ze_0$
and $\mathbb Z/2$.  Every nonidentity power of each specified affine
action $\tau_i$ is free.
\end{theorem}

\begin{theorem}\label{night:P1.3}
\StatementRecord{night:P1.3}{proved}{}{Definitions and the displayed constructions}
The two finite filling pieces have first integral homology
$\mathbb Z^2$ and $\mathbb Z^2\oplus\mathbb Z/2$.
\end{theorem}

\begin{lemma}[The actual order-four torsion transgression]
\label{mapsaudit:order-four-torsion-transgression}
\StatementRecord{mapsaudit:order-four-torsion-transgression}{proved}{}{\cref{maps:src:F.2.5}; the integral quotient and explicit certificates in \cref{maps:src:F.2.6}; \cref{maps:eq:F.6}; the displayed cellular and boundary maps}

Use the order-four filling and the marked matrices of \cref{maps:src:F.2.5}.
Let $T=\bigwedge^{q-1}D^t$ on the integral vertical cochains of the
five-torus. If $w$ represents a torsion element of
$\operatorname{coker}(T-I)$, its point-supported degree-$q$ class has
Leray differential equal, up to the one boundary-orientation sign, to
the image of $w$ in the global degree-$(q-1)$ coinvariants. There is
no multiplicity-four factor. For each of $q=2,3,4$, the unique nonzero local torsion class maps
to the nonzero class in the global $\mathbb Z/2$ quotient.
\end{lemma}

\begin{theorem}\label{night:P1.10}
\StatementRecord{night:P1.10}{proved}{}{\cref{night:P1.11}, \cref{night:P1.9}}
For every real parameter $c>-h_{\min}$ the corrected construction gives
a compact smooth complex fourfold $M_c$ and a proper holomorphic map
to $\mathbb P^1$ with the two specified multiple fibres and the
index-one Kuhn cusp.  With its complex orientation,
$\pi_1(M_c)=1$ and $\chi(M_c)=6$.
\end{theorem}

\begin{theorem}\label{night:P1.11}
\StatementRecord{night:P1.11}{proved}{}{\cref{night:P1.9}}
The index-one central cusp fibre $C$ has torsion-free integral homology
with ranks $(1,3,9,9,10,3,1)$.
The collapse $\kappa:T^6\to C$ has primitive image in every degree;
it is onto except in degree four, where its cokernel is $\mathbb Z$.
Its image ranks are $(1,3,9,9,9,3,1)$.
In every degree its kernel is the saturation of
$(\bigwedge^kM_0-I)\bigwedge^k\Lambda$.
\end{theorem}

\begin{theorem}\label{night:P1.12}
\StatementRecord{night:P1.12}{proved}{}{\cref{night:P1.1}, \cref{night:P1.9}, \cref{night:P1.10}, \cref{night:P1.11}; the actual marked evaluation and finite logarithmic collar maps; \cite[Theorem 4.1.1 and Lemmas 4.1.2--4.1.3, pp.~480--482]{ffAchingerOgus20}; \cite[Theorem 3.2, p.~195; Theorem 3.30, p.~241; Theorem 4.32, p.~366]{Hatcher2002}}
The fourfold $M_c$ has torsion-free integral homology
\[
H_k(M_c;\mathbb Z)=
\begin{cases}
\mathbb Z,&k=0,8,\\
\mathbb Z^4,&k=4,\\
0,&k=1,2,3,5,6,7.
\end{cases}
\]
In particular it is three-connected.
\end{theorem}

\begin{theorem}\label{night:P1.19}
\StatementRecord{night:P1.19}{proved}{}{\cref{night:P1.11}, \cref{night:P1.12}, \cref{night:P1.9}}
There is an independent computation from the full global cellular
complex.  Its chain ranks in degrees zero through eight are
\[
 (4,26,87,162,200,156,79,23,3),
\]
and its differentials, after the explicit integral chain equivalences
below, are the matrices of case one in
\path{code/part1/global_ext_cases.json}.
Their integral Smith calculation gives
\[
 H_k(M;\mathbb Z)=
 \begin{cases}
 \mathbb Z,&k=0,8,\\
 \mathbb Z^4,&k=4,\\
 0,&\text{otherwise}.
 \end{cases}
\]
This calculation uses local Poincar\'e--Lefschetz duality of the
smooth cusp cap to determine its secondary attaching class.  It does
not use global Poincar\'e duality or the homology conclusion of P1.12.
In particular, the map obtained by setting the whole cusp suspension
chain term to zero is not the actual relative cusp model; its precise
missing term is accounted for here.
\end{theorem}

\begin{theorem}\label{night:P1.4}
\StatementRecord{night:P1.4}{proved}{}{Definitions and the displayed constructions}
The genus-two pencil has one nodal fibre away from $0,\infty$, at
$\lambda_*=-256/3125$.  Good reduction at $0,\infty$ requires root
orders $8,5$.  On this root orbifold its Hodge bundle is
$\mathcal H=T^3\oplus T$, where
\[
T=\mathcal O\bigl((2/5)[\infty]-(3/8)[0]\bigr),\qquad
\deg T=1/40,
\]
and $T^3$ is the line generated meromorphically by $dx/y$.
\end{theorem}

\begin{theorem}\label{night:P1.5}
\StatementRecord{night:P1.5}{proved}{}{Definitions and the displayed constructions}
The formula in the prompt with $P$ the image of
$\ell_0\wedge\mathcal H$, and $\mathcal H$ the rank-two Hodge bundle,
cannot produce a rank-three subbundle.  The required rank-three image
is instead the image of $\ell_0\wedge\mathcal H_{\mathrm{dR}}$, where
$\mathcal H_{\mathrm{dR}}$ is the full rank-four de Rham bundle.
\end{theorem}

\begin{theorem}\label{night:P1.2}
\StatementRecord{night:P1.2}{proved}{}{Definitions and the displayed constructions}
The lower blocks of $A,B$ are precisely the quotient representations of
$\bigwedge^2(C_f^{-1}),\bigwedge^2C_g$ in the supplied basis $Q$.
The integral quotient is the dual of the primitive exterior-square
lattice.  The natural map from the primitive lattice into that quotient
has cokernel $\mathbb Z/2$; these two integral representations are not
isomorphic.  Their rationalizations are isomorphic by that same map.
\end{theorem}

\begin{theorem}\label{night:P1.7}
\StatementRecord{night:P1.7}{proved}{}{\cref{maps:src:F.3.5}, \cref{night:P1.2}}
The integral $H^1$ monodromy of the pencil has a companion marking with
finite pair $(C_f^{-1},C_g)$, and cusp product $C_fC_g^{-1}$.
Its exterior-square quotient is consequently the integral quotient
representation appearing in P1.2.
\end{theorem}

\begin{theorem}\label{night:P1.8}
\StatementRecord{night:P1.8}{proved}{}{\cref{night:P1.5}}
With the corrected rank-three bundle $P$ of P1.5, the supplied rank-six
system admits a holomorphic lifted family
$F_c=(\sigma+ic\alpha e)(P)$ on the punctured base.
For a specified splitting $\sigma$ there is a real analytic function
$h$ with finite attained minimum $h_{\min}$, and the exact domain of
real parameters giving $F_c\cap\overline{F_c}=0$ everywhere is
$c>-h_{\min}$.
At the cusp its periods have the form
$Z(q)=(\log q/(2\pi i))I_3+C(q)$ with $C$ holomorphic and bounded.
The construction extends to the two finite root charts and descends to
orders $5,4$ on the rank-six system.
\end{theorem}

\begin{theorem}\label{night:P1.9}
\StatementRecord{night:P1.9}{proved}{}{Definitions and the displayed constructions}
The periodic Kuhn fan, in the marking $S$, gives a smooth proper
holomorphic filling of the periods $Z(q)=sI_3+C(q)$ for a sufficiently
small disk.  Its central fibre is reduced with normal crossings.
Modulo the index-one period lattice it has respectively $1,7,12,6$
strata of dimensions $3,2,1,0$, each isomorphic to the corresponding
complex torus $(\mathbb C^*)^d$, and Euler number $6$.
An index-$d$ period sublattice gives $d$ times these stratum counts and
Euler number $6d$.
\end{theorem}

\begin{theorem}\label{night:P1.13}
\StatementRecord{night:P1.13}{proved}{}{\cref{night:P1.12}; \S\ref{sec:V-known-inputs}; \cite[equations (1.2)--(1.4), pp.~3--4]{ffEsoleKang}; \cite[Theorem 2.2.1, p.~61]{ffPetersSterk}}
With the complex orientation the fourfold has
\[
\sigma(M_c)=2,\qquad H^4(M_c;\mathbb Z)\cong I_{3,1},\qquad
w_2(M_c)=0,
\]
and there is a diagonal integral basis in which
$p_1/2=(1,1,1,1)$.
Its Chern numbers are
$c_4=6$, $c_2^2=2$, and
$c_1^4=c_1^2c_2=c_1c_3=0$; its Pontryagin numbers are
$p_1^2=8$, $p_2=14$.
\end{theorem}

\begin{theorem}[a complex structure on the stated connected sum]\label{maps:src:F.4.1}
\StatementRecord{maps:src:F.4.1}{proved}{}{the hypotheses, notation and definitions stated here; \cref{maps:src:F.3.7}; \cref{maps:src:F.2.6}; \S\ref{sec:V-known-inputs}; the explicit integral reflections below; \cite[Corollary 1.4, p.~4]{ffMandal}; \cite[Definition 3.6 and Theorem 3.7(1), pp.~12--14]{IVVWallForm}; \cite[Theorem A, p.~178]{IVVKramerStolz}; \cref{night:P1.13}}


The compact complex fourfold of Theorem \ref{maps:src:F.3.7} is three-connected,
has Euler number six, and is orientation-preservingly diffeomorphic to
\[
 3\mathbb H\mathbb P^2\ \#\ (-\mathbb H\mathbb P^2).
 \tag{F.25}\label{maps:eq:F.25}
\]
Its integral middle lattice is $I_{3,1}$ and, in a diagonal integral
basis, $p_1/2=(1,1,1,1)$.  In particular the conclusion is an exact
smooth diffeomorphism, with no remaining homotopy-sphere ambiguity.

\end{theorem}

\begin{theorem}\label{night:P1.15}
\StatementRecord{night:P1.15}{proved}{}{\cref{night:P1.1}, \cref{night:P1.8}}
The rank-six flat extension, as a complex flat bundle with its invariant
line, has the following explicit algebraic connection model.  Set
$\omega_k=x^kdx/y$, $0\leq k\leq3$, and use the quotient frame
\[
(w_2,w_1,w_0,w_{-1},w_{-2})
 =(\omega_{01},\omega_{02},\omega_{03},\omega_{13},\omega_{23}),
\qquad \omega_{12}=-3\omega_{03}.
\]
These are meromorphic frames of
$(L^2,L,\mathcal O,L^{-1},L^{-2})$, where $L=T^2$ on the root
orbifold of orders $5,4$.  If a frame is written as a row of vectors,
its derivative is right multiplication by
\[
\Gamma_6=\begin{pmatrix}0&b\\0&\Gamma_W\end{pmatrix},\qquad
b=\frac1{\lambda-\lambda_*}
 \left(-\frac{25}{6\lambda^2},-\frac{125}{48\lambda},0,1,0\right),
\tag{1.15.1}
\]
where
\[
\resizebox{\textwidth}{!}{$\displaystyle
\Gamma_W=\frac1{3125\lambda+256}
\begin{pmatrix}
-8(625\lambda+48)/\lambda&-8&10\lambda&75\lambda^2/2&0\\
-125/(2\lambda)&-4(625\lambda+48)/\lambda&8&0&75\lambda^2/2\\
300/\lambda^2&-375/\lambda&0&-48&60\lambda\\
120/\lambda^3&0&125/(2\lambda)&4(625\lambda+48)/\lambda&-8\\
0&120/\lambda^3&50/\lambda^2&-125/(2\lambda)&8(625\lambda+48)/\lambda
\end{pmatrix}.$}
\tag{1.15.2}
\]
This model specifies a canonical algebraic splitting, denoted
$\sigma_b$.  The compatible contraction used for the real parameter is
$\alpha=(625\pi i/128)\alpha_{\rm alg}$, where
$\alpha_{\rm alg}(w_0)=4/3$ and it vanishes on $w_1,w_2$.
\end{theorem}

\begin{theorem}\label{night:P1.16}
\StatementRecord{night:P1.16}{proved}{}{\cref{night:P1.15}, \cref{night:P1.18}, \cref{night:P1.8}}
In the canonical splitting $\sigma_b$ of P1.15,
\[
\begin{aligned}
-1.0342766224163217994942488453783461&<h_b(1),\\
h_b(1)&<-1.0342766224163217994942488453783460.
\end{aligned}
\]
In particular $h_b(1)<-1$.  The parameter $c=5/(3\pi)$ is outside the
global transversality domain.  The lifted period matrices at
$\lambda=1,2$ have been evaluated by two different continuation methods
agreeing to more than thirty decimal digits.
\end{theorem}

\begin{theorem}\label{night:P1.18}
\StatementRecord{night:P1.18}{proved}{}{\cref{night:P1.15}; the rational connection and root-gauge definitions displayed in \cref{night:P1.16}; \cref{night:P1.8}; \cite[Lemma~2.7, pp.~42--43]{TeschlODE2012}; \cite[\S1.9(iii), (1.9.30)--(1.9.31)]{DLMFCauchy2026}}
In the algebraic splitting $\sigma_b$, the function $h_b$ has its
unique global minimum at the finite-root point $\lambda=\infty$.
The minimum is the exact Frobenius period constant
\[
 h_{b,\min}=-\frac{96}{625\pi}\operatorname{Re}(XH_\infty)_{03},
 \qquad
 h_{b,\min}=-1.0945402503114027801162682630549457037\ldots,
\tag{1.18.1}
\]
where $X$ is the normalized integral intertwiner of P1.15 and
$H_\infty$ is the regular infinity-root limit of its fundamental
solution, normalized to be the identity at $\lambda=1$ before
multiplication by $X$.  Formula (1.18.1), the rational connection,
and those two normalizations specify the constant exactly.
A fresh certified enclosure has radius less than $4\cdot10^{-76}$
about the full stored midpoint in
\path{code/paper_v/output/fresh_frobenius.json}.
Consequently the exact domain of completed transverse periods is
\[
 c>-h_{b,\min},\qquad
 \operatorname{Im}\xi>-h_{b,\min}
\]
for the real and complex parameters, respectively.
The member $c=10$ used in the numerical period checks is therefore
admissible on the whole completed base.

For the point-normalized splitting of P1.8, continued from
$1$ to $\lambda_*/2$ by the polygon
$1,1+i,\lambda_*/2+i,\lambda_*/2$, the exact comparison constant
$a_{\rm split}$ of (1.15.4) has certified value
\[
\begin{split}
 \operatorname{Re}a_{\rm split}
   &=-0.7099885679126427771723473455012308149\ldots,\\
 \operatorname{Im}a_{\rm split}
   &=-0.9134017059364153362143229223828373269\ldots.
\end{split}
\]
Thus the minimum in that splitting is
$-0.1811385443749874439019453406721083768\ldots$,
attained at the same base point.  Its real parameter interval is
\[
 c>0.1811385443749874439019453406721083768\ldots.
\]
The full parameter comparison remains
$\xi_b=\xi_{\rm original}-a_{\rm split}$; both its real and imaginary
parts are necessary to identify the two normalized families.
\end{theorem}

\begin{theorem}\label{night:P1A.1}
\StatementRecord{night:P1A.1}{proved}{}{\cref{night:P1.10}, \cref{night:P1.8}, \cref{night:P1.9}}
For every admissible member, its vertical translation bundle on
$\widehat B$ and its first coherent direct image are
\[
 E_0=F^\vee=\mathcal O\oplus L^{-1}\oplus L^{-2},\qquad
 G=R^1h_*\mathcal O_M=\mathcal O\oplus L^{-1}\oplus L^{-2}.
\]
Moreover $R^qh_*\mathcal O_M=\bigwedge^qG$ for $0\le q\le3$,
$K_M=h^*L^{-6}$, and
\[
 (h^0(\mathcal O_M),\ldots,h^4(\mathcal O_M))=(1,1,0,0,0).
\]
\end{theorem}

\begin{theorem}\label{night:P1A.2}
\StatementRecord{night:P1A.2}{proved}{}{\cref{night:P1.15}, \cref{night:P1A.1}}
The complete Hodge table of every admissible fourfold is the
following, with rows indexed by $p$ and columns by $q$:
\[
\begin{array}{c|ccccc}
h^{p,q}&0&1&2&3&4\\\hline
0&1&1&0&0&0\\
1&0&2&1&0&0\\
2&0&1&6&1&0\\
3&0&0&1&2&0\\
4&0&0&0&1&1
\end{array}
\]
\end{theorem}

\begin{theorem}\label{night:P1A.3}
\StatementRecord{night:P1A.3}{proved}{}{\cref{night:P1.9}, \cref{night:P1.10}, \cref{night:P1.15}, \cref{night:P1A.1}, \cref{night:P1A.2}}
For every admissible fourfold,
\[
 (h^0(T_M),h^1(T_M),h^2(T_M),h^3(T_M),h^4(T_M))
 =(1,1,1,2,0).
\]
The complex parameter family has a nonzero Kodaira--Spencer map
at every member and supplies a smooth one-dimensional miniversal
Kuranishi germ.
\end{theorem}

\begin{theorem}\label{night:P1A.4}
\StatementRecord{night:P1A.4}{proved}{}{\cref{night:P1.1}, \cref{night:P1.9}, \cref{night:P1.15}, \cref{night:P1A.1}}
The fourfold has algebraic dimension one and
\[
\Aut^0(M)=\mathbb C,\qquad
\Aut(M)=\mathbb C\times\mathbb Z/2.
\]
In a fixed holomorphic splitting, two members of the full
admissible complex parameter half-plane are biholomorphic
exactly when their parameters are equal.
\end{theorem}

\begin{theorem}\label{night:P1.17}
\StatementRecord{night:P1.17}{proved}{}{\cref{night:P1.1}}
For the cusp marking $S$ fixed in P1.1 and the corresponding Kuhn fan,
none of the completed fourfolds in this family admits an
anti-holomorphic automorphism.  In particular none admits a real
structure.  More strongly, two members with this same cusp marking
are never anti-biholomorphic.
\end{theorem}

\begin{theorem}\label{night:P1A.5}
\StatementRecord{night:P1A.5}{proved}{}{\SeriesRef{night:P2.4}, \cref{night:P1A.1}, \cref{night:P1.8}; \cite[Theorem 1.6, preprint p.~5]{BM97}; \cite[\S6, Hauptsatz I; \S7, Satz 1 and Satz 5]{Grauert60}}
There is no such analytic extension across any finite boundary
point of $\mathcal H$, or across infinity. The same conclusion
holds after any finite ramified change of the local boundary
coordinate. In particular neither end of the real interval
$c_*<c<+\infty$ has a proper Hausdorff analytic extension in
this sense.
\end{theorem}

\begin{theorem}\label{night:P1A.6}
\StatementRecord{night:P1A.6}{proved}{}{\cref{night:P1A.3}, \cref{night:P1A.4}}
A holomorphic punctured-disc family whose fibres are all
members of this fourfold family has an interior parameter
limit. After a base change of degree at most two, it is
isomorphic to the pullback of the constructed family by
a holomorphic map of the full disc to $\mathcal H$, and
therefore has a proper smooth filling. Such a family
cannot approach either parameter end.
\end{theorem}

\begin{theorem}\label{night:P1B.1}
\StatementRecord{night:P1B.1}{proved}{}{Definitions and the displayed constructions}
Let $N$ be an endomorphism of $\Z^6$ with $N^2=0$, rank three,
and primitive image. There is a unimodular basis in which
\[
 I+N=\begin{pmatrix}I_3&I_3\\0&I_3\end{pmatrix}.
\]
Its centralizer in $\operatorname{GL}_6(\Z)$ consists exactly of
$\left(\begin{smallmatrix}P&Q\\0&P\end{smallmatrix}\right)$,
with $P\in\operatorname{GL}_3(\Z)$ and $Q\in\operatorname{Mat}_3(\Z)$.
\end{theorem}

\begin{theorem}\label{night:P1B.2}
\StatementRecord{night:P1B.2}{proved}{}{Definitions and the displayed constructions}
A finite-order integral linear action on $\R^6$ admitting a
free finite affine torus lift has characteristic polynomial
among exactly thirty-eight cyclotomic products. Their linear
orders are
\[
 1,2,3,4,5,6,8,10,12.
\]
Exactly nineteen of these polynomials admit a commuting complex
structure on $\R^6$. These are a preliminary spectral list;
no assertion about global gluings is inferred from the list.
\end{theorem}